\documentclass[11pt]{article}
\usepackage[margin=1in]{geometry}
\usepackage[pdftex]{graphicx}
\usepackage{amsmath,amssymb,amsthm}
% Latin Modern is Computer Modern redrawn as a scalable outline font. custom.sty
% switches on microtype's automatic font expansion, which pdfTeX can only do
% with scalable fonts -- and the bold dollar signs in the fees section pull in a
% bitmap text-companion face that aborts the build. Loading lmodern first makes
% every face scalable; the document looks the same as the other handouts.
\usepackage{lmodern}
\usepackage[T1]{fontenc}
% The timeline runs longer than the space left on its page. A plain tabular
% cannot break, so it would be shunted whole onto the next page and leave a
% third of a page blank; longtable lets it flow.
\usepackage{longtable}
\usepackage{custom}

% color boxes
\usepackage[many]{tcolorbox}
\tcbuselibrary{theorems}
\newtcolorbox{official}{colback=red!8,colframe=red!55!black,breakable,enhanced,
  left=8pt,right=8pt,top=6pt,bottom=6pt}
\newtcolorbox{keybox}{colback=blue!8,colframe=blue!40!black,breakable,enhanced,
  left=8pt,right=8pt,top=6pt,bottom=6pt}
\newtcolorbox{tipbox}{colback=green!8,colframe=green!40!black,breakable,enhanced,
  left=8pt,right=8pt,top=6pt,bottom=6pt}

% headers and footers
\usepackage{fancyhdr}
\pagestyle{fancy}
\lhead{F = ma Exam FAQ}
\chead{}
\rhead{EichenlaubPhysics}
\lfoot{}
\cfoot{\thepage}
\rfoot{}
\renewcommand{\headrulewidth}{0.4pt}
\setlength{\headheight}{14pt}

\newcommand{\psettitle}[1]{
    \begin{center}
    \huge \textbf{#1}
    \end{center}
}

% A FAQ question. Kept with at least the start of its answer so a question never
% ends up stranded alone at the foot of a page.
\definecolor{qblue}{rgb}{0.10,0.22,0.50}
\newcommand{\Q}[1]{\par\medskip\noindent\textbf{\color{qblue}#1}\par\nobreak\vspace{1pt}\nobreak}

\linespread{1.03}
\setlength{\parindent}{0pt}
\setlength{\parskip}{4pt}

\begin{document}

\psettitle{F = ma Exam FAQ}
\begin{center}
\large Everything about \emph{taking} the exam --- the logistics, not the physics
\end{center}

\begin{official}
\textbf{The AAPT website is the final word.}\\[4pt]
Everything in this handout is my summary, written to save you time. If anything here
disagrees with AAPT, \textbf{AAPT is right and I am wrong}. Check the official site
before you rely on any date, fee, or rule:
\begin{center}
\large\bfseries
\href{https://www.aapt.org/physicsteam/}{www.aapt.org/physicsteam}
\end{center}
The two pages worth bookmarking are
\href{https://www.aapt.org/physicsteam/registration.cfm}{Registration} and
\href{https://www.aapt.org/physicsteam/PT-rules.cfm}{Rules}.

\textbf{A note on dates and fees.} The numbers below are from the \textbf{2026 cycle},
which is the most recent one AAPT has published in full. At the time I wrote this, AAPT
had not yet posted the next cycle's page. Use the 2026 figures to understand the
\emph{shape} of the year --- roughly when things happen and roughly what they cost ---
and confirm the actual dates on the AAPT site once they are announced. Dates move.
\end{official}

\section{The Short Version}

\begin{keybox}
\textbf{The F = ma exam is 25 multiple-choice questions in 75 minutes, on mechanics
only.} You take it \textbf{in person, on a computer, with a live proctor watching}, on a
single day in February. \textbf{Your school has to register you} --- you cannot sign
yourself up --- and registration closes about three weeks before the exam. Score high
enough (recently 14--18 out of 25) and you are invited to the USAPhO exam in April.
\end{keybox}

\textbf{The single most important thing in this document:} the step most likely to stop
you is not the physics. It is finding a teacher willing to proctor and getting your school
registered before the deadline. Start asking in the fall.

\section{Timeline: What To Do, and When}

Dates are the 2026 cycle's; treat the left column as your real deadline structure.

\begin{longtable}{p{0.24\textwidth}p{0.66\textwidth}}
\hline
\textbf{When} & \textbf{What you should be doing} \\
\hline
\endfirsthead
\hline
\textbf{When} & \textbf{What you should be doing} \emph{(continued)} \\
\hline
\endhead
September--October &
Ask your physics teacher whether the school will run the F = ma exam. If they have never
heard of it, send them the AAPT link. This is the step that fails most often, and it fails
because people leave it until January. \\[3pt]
October--November &
Confirm \emph{who} is registering and whether the school will pay the fee. Get a name, not
a vague yes. If your school says no, start looking elsewhere now (see below). \\[3pt]
December--early January &
Make sure the proctor has actually registered. Do not assume. Ask them to confirm they
have submitted it. \\[3pt]
Mid-January &
\textbf{Registration deadline.} In 2026 this was January 20 (extended to January 21). Once
it passes, there is no way in for that year. \\[3pt]
Mid-February &
\textbf{Exam day.} In 2026 it was February 12, in a 1:00--4:00 pm EST window --- your
proctor picks the start time inside that window. In 2024 and 2023 it was February 8 and 9. \\[3pt]
Early March &
Solutions posted (March 5 in 2026), then results (March 9 in 2026). \\[3pt]
April &
\textbf{USAPhO}, if you qualified. April 10 in 2026, a 1:00--5:00 pm EDT window. \\[3pt]
May--June &
USAPhO Plus, then the roughly 20 students invited to the training camp
(May 30--June 9, 2026, at the University of Maryland). \\
\hline
\end{longtable}

\section{The Exam Itself}

\Q{What is on it?}
Mechanics only --- kinematics, Newton's laws, energy, momentum, rotation, oscillations,
gravitation, fluids. No electricity and magnetism, no thermodynamics, no optics, no modern
physics.

\Q{How long is it and how many questions?}
25 multiple-choice questions in 75 minutes. That is three minutes per question on average.

\Q{Do I need calculus?}
No. Every question can be done without it. Calculus occasionally makes a problem faster,
but nothing on the exam requires it.

\Q{How is it scored? Is there a guessing penalty?}
One point for each correct answer, and \textbf{zero points for a wrong answer or a blank}.
AAPT's own instructions say ``There is no additional penalty for incorrect answers.''

\begin{tipbox}
\textbf{So never leave a question blank.} A blank and a wrong answer score identically, and
a guess among five choices is worth $0.2$ points on average. If you have eliminated even one
option, guess. With two minutes left, put an answer on everything you have not reached.
\end{tipbox}

\Q{Are the questions in order of difficulty?}
Roughly, but not reliably. AAPT states that ``All questions are equally weighted, but are
not necessarily equally difficult.'' Since every question is worth the same one point,
there is no reason to grind on question 8 while questions 20 and 21 --- which you might find
easier --- sit unread. Take a pass through the whole exam.

\Q{What value of $g$ should I use?}
The exam tells you: use $g = 10$ N/kg unless a question says otherwise. This is deliberate
--- it keeps the arithmetic clean.

\section{How You Actually Take It}

\Q{Is it on paper or on a computer?}
On a computer. As of the 2026 cycle, AAPT states that the F = ma and USAPhO exams are
``conducted online via Educational Vistas only,'' and that ``There will be no printable
versions of the exams available.'' This changed recently --- older exams were paper with a
bubble sheet, and in 2021 the platform was Art of Problem Solving --- so ignore what you
read in older guides.

\Q{Can I take it at home?}
No. It is proctored and taken live, in person, at your school or another approved site,
with everyone at that site testing at the same time. It is not a take-home or an
open-window online test.

\Q{Do I get scratch paper?}
Yes, and only from the proctor. The rule is that ``Only scratch paper provided by the
proctors may be used,'' and ``The proctors must collect it immediately after the exam.''
So you cannot bring your own paper, and you do not get your work back.

\Q{What are the calculator rules?}
You may use a hand-held calculator, subject to all of the following:

\begin{itemize}\setlength{\itemsep}{1pt}
\item \textbf{Not a graphing calculator.}
\item Its display has \textbf{no more than three lines}.
\item Its \textbf{memory must be completely cleared} of data and programs immediately
before the exam.
\item \textbf{Only the basic functions} found on a simple scientific calculator.
\item \textbf{Calculators may not be shared.}
\end{itemize}

A cheap scientific calculator is the safe choice. Do not show up with a TI-84 and expect
to use it. Bring your own, and clear it before you sit down.

\Q{What else is banned?}
Cell phones may not be used during the exam or while the exam materials are present, and
you may not use any external references --- no books, no formula sheets, no notes. There is
no formula sheet provided either.

\Q{Can I talk about the questions afterwards?}
Not right away. Each exam carries an embargo: the 2025 paper told students not to
communicate anything about the questions or their solutions until after February 13, 2025 ---
the day after the exam. Respect it. Solutions are posted publicly a few weeks later.

\section{Getting a Seat: Proctors, Registration, and Fees}

\Q{Who signs me up?}
Not you. \textbf{Your school registers you.} AAPT says ``It is preferred that one person
register all students from a school,'' and that person is normally the teacher who will
proctor. Your job is to find that person and make sure they do it before the deadline.

\Q{Who is allowed to proctor?}
\begin{itemize}\setlength{\itemsep}{1pt}
\item A proctor ``should have at least a 2-year degree, though this degree does not need to
be related to physics.'' Your English teacher can proctor.
\item \textbf{Not} a parent, relative, or close friend or acquaintance of a student taking
the exam.
\item \textbf{Not} a for-profit test prep centre or educational service --- AAPT
specifically disallows this now.
\item They should register with an official school email address where possible.
\end{itemize}

\Q{What does it cost?}
In 2026: \textbf{\$75 per school} (\$37.50 if the registering teacher is an AAPT member)
\textbf{plus \$15.50 per student.} Note the shape of that --- most of the cost is a flat
per-school fee, so the more students who sign up, the cheaper it gets per person.

\Q{Do I have to pay it?}
Often not. \textbf{Many schools cover the fee} out of a department or activities budget,
especially if a few students sign up together. Ask --- do not assume you will be billed. If
cost is genuinely the obstacle, tell me and we will sort it out.

\begin{tipbox}
\textbf{Recruit a friend or two.} It makes the per-student cost lower, it makes the school
far more likely to say yes, and it gives the teacher a reason to run it. ``Can you proctor
an exam for five of us?'' is a much easier ask than ``Can you proctor an exam for me?''
\end{tipbox}

\Q{My school says no, or has no physics teacher. Now what?}
You are not stuck. You need a proctor and a site, and it does not have to be your own
school:

\begin{itemize}\setlength{\itemsep}{1pt}
\item \textbf{Another school.} A nearby high school already running the exam can usually
add one more student. This is common and entirely allowed.
\item \textbf{A community college or university.} Many have a testing centre.
\item \textbf{A public library.} Some offer proctoring services.
\end{itemize}

AAPT's advice to homeschooled students is exactly this --- contact community colleges,
universities, or libraries. The same route works for anyone whose school will not host it.

\Q{Will AAPT find me a proctor or a test centre?}
No. AAPT is explicit that it ``does not facilitate or arrange test centers, teachers, or
proctors,'' and does not keep a list you can consult. Finding a site is on you. Start early
and ask several places at once.

\Q{Am I eligible?}
You must be a U.S. citizen, a U.S. permanent resident (green card holder), or currently
attending a U.S. school --- \textbf{and you must be physically located in the U.S. when you
take it}. There is no grade-level or age restriction on the F = ma exam itself. Age limits
apply only later, to the training camp and the international team.

\Q{Can I take it more than once in a year?}
Not currently. There is one exam per year. In 2022 there were two (an A and a B version)
and students could sit both, but every year since has had a single exam. If you miss it,
you wait a year.

\section{What Score Do I Need?}

\Q{What is the cutoff?}
There is no fixed cutoff. AAPT sets it \emph{after} the exam, at whatever score invites
roughly 400--500 students to the USAPhO. So it moves every year with the difficulty of the
paper. Here is the actual history:

\begin{center}
\begin{tabular}{llrrrcl}
\hline
\textbf{Year} & \textbf{Exam date} & \textbf{Took it} & \textbf{Median} & \textbf{Mean} &
\textbf{Cutoff} & \textbf{Qualified} \\
\hline
2021       & Feb 18 & 3524 & 8  & 8.89  & 15 & --- \\
2022 (A)   & Feb 9  & 2038 & 8  & 8.24  & 15 & --- \\
2022 (B)   & Feb 15 & 2518 & 8  & 8.51  & 15 & --- \\
2023       & Feb 9  & 5173 & 10 & 10.29 & 18 & 417 \; (8.1\%) \\
2024       & Feb 8  & 6061 & 7  & 7.93  & 14 & 500 \; (8.2\%) \\
2025       & Feb 12 & 5848 & 7  & 8.21  & 15 & 518 \; (8.9\%) \\
2026       & Feb 12 & 6682 & 8  & 8.90  & 16 & 440 \; (6.6\%) \\
\hline
\end{tabular}
\end{center}

\emph{Source: AAPT's official ``F = ma Contest Results'' documents for each year.}

\Q{So what should I actually aim for?}
Read the table two ways.

First, the cutoff has sat between \textbf{14 and 18} for six years running. Aiming at 14
because ``that was enough in 2024'' is how people miss by a point --- in 2023 the same
performance would have been four points short. \textbf{Target 18--20.} That clears every
recent cutoff with room for a bad day.

Second, look at the medians: 7 to 10 out of 25. \textbf{Half of everyone who sits this
exam scores single digits}, and these are self-selected students who chose to take a
physics olympiad qualifier. Qualifying means finishing in roughly the top 7--9\% of that
group. A score you might think of as mediocre --- 12, 13 --- is well above the national
median. Do not measure yourself against 25.

\Q{Does my score get published?}
No. Individual scores are not disclosed publicly. AAPT publishes the aggregate statistics
and the cutoff, and notifies qualifiers.

\section{After the Exam}

\Q{What happens if I qualify?}
You are invited to the \textbf{USAPhO} --- a much longer free-response exam, also
proctored, held in April (April 10 in 2026, scheduled in a 1:00--5:00 pm EDT window). It is
a completely different animal: a handful of long multi-part problems where you show all
your work, not multiple choice. Check AAPT for the exact length and structure in the year
you sit it.

\Q{And after that?}
Top USAPhO scorers take the USAPhO Plus exam, from which about 20 students are invited to
the training camp at the University of Maryland in late May. Five are eventually selected
for the U.S. team at the International Physics Olympiad.

\Q{What if I do not qualify?}
Almost nobody qualifies the first time they sit it. The exam is offered every year, the
scores in the table above tell you how selective it is, and the ordinary path is to take
it, see where you stand, and come back stronger. That is what the practice tests and the
handouts are for.

\section{Your Checklist}

\begin{keybox}
\begin{enumerate}\setlength{\itemsep}{2pt}
\item[$\square$] Asked a teacher about proctoring \textbf{(do this in the fall)}
\item[$\square$] Confirmed who is registering and whether the school pays
\item[$\square$] Recruited a friend or two to bring the cost down
\item[$\square$] Confirmed the proctor actually submitted the registration
\item[$\square$] Checked the exam date and start time with the proctor
\item[$\square$] Bought a non-graphing scientific calculator and cleared its memory
\item[$\square$] Know that blanks score zero, so nothing gets left blank
\end{enumerate}
\end{keybox}

\vspace{0.5em}
\begin{center}
\emph{Questions this handout does not answer? Ask me --- and check
\href{https://www.aapt.org/physicsteam/}{aapt.org/physicsteam}, which is always the final
authority.}
\end{center}

\end{document}
